% Layout fixture: lecture slides (beamer class, default theme). Written for ReadableSTEM's
% tests; no third-party text.
\documentclass{beamer}
\usepackage{amsmath}
\usepackage{tikz}
\setbeamertemplate{navigation symbols}{}

\begin{document}

\begin{frame}{Motion in One Dimension}
  \begin{itemize}
    \item Position $x$ is measured from a chosen origin.
    \item Velocity is the rate of change of position:
    \begin{itemize}
      \item average velocity uses two times,
      \item instantaneous velocity uses a limit.
    \end{itemize}
    \item Acceleration is the rate of change of velocity.
  \end{itemize}
\end{frame}

\begin{frame}{Constant Acceleration}
  When the acceleration $a$ is constant,
  \begin{align*}
    v &= v_0 + a t, \\
    x &= x_0 + v_0 t + \tfrac{1}{2} a t^2 .
  \end{align*}
  Eliminating the time gives
  \begin{equation*}
    v^2 = v_0^2 + 2 a \left(x - x_0\right).
  \end{equation*}
\end{frame}

\begin{frame}{Velocity Graph}
  \begin{figure}
    \centering
    \begin{tikzpicture}[scale=0.9]
      \draw[->] (0,0) -- (5,0) node[right] {$t$};
      \draw[->] (0,0) -- (0,3) node[above] {$v$};
      \draw[thick] (0,0.5) -- (4.5,2.6);
    \end{tikzpicture}
    \caption{Velocity rises linearly when the acceleration is constant.}
  \end{figure}
\end{frame}

\end{document}
